\documentclass{article}

\usepackage[spanish]{babel}
\usepackage{amsmath,amssymb,amsthm}

\newcommand{\N}{\mathbb{N}}
\newcommand{\Z}{\mathbb{Z}}
\newcommand{\Q}{\mathbb{Q}}
\newcommand{\R}{\mathbb{R}}
\newcommand{\C}{\mathbb{C}}

\newcommand{\vect}[1]{\vec{#1}}
\newcommand{\inner}[2]{\langle #1, #2 \rangle}
\newcommand{\norm}[1]{\| #1 \|}

\theoremstyle{definition}
\newtheorem{definition}{Definición}[section]
\newtheorem{example}[definition]{Ejemplo}

\theoremstyle{plain}
\newtheorem{theorem}[definition]{Teorema}
\newtheorem{proposition}[definition]{Proposición}
\newtheorem{lemma}[definition]{Lema}

\theoremstyle{remark}
\newtheorem{remark}[definition]{Observación}

\title{Calculo en varias variables}
\author{Benjamín Rivas Beltrán}
\date{\today}

\begin{document}

\maketitle

\section{Nociones básicas de $\R^n$}

$(\R^n, +, \cdot)$ es un espacio vectorial  sobre $\R$.

\begin{definition}[Producto interior en $\R^n$]

Sean $\vect{x}, \vect{y} \in \R^n$ con:
\[
\vect{x} = (x_1, x_2, \dots, x_n) \quad \text{y} \quad \vect{y} = (y_1, y_2, \dots, y_n) \]
El producto interior de $\vect{x}$ y $\vect{y}$ es el número real dado por:
\begin{align*}
\inner{\vect{x}}{\vect{y}} &= x_1 y_1 + x_2 y_2 + \dots + x_n y_n \\
&= \sum_{i=1}^n x_i y_i
\end{align*}

\end{definition}

\begin{example}

Sea $\vect{x} = (1, 2, 3)$ y $\vect{y} = (4, 5, 6)$. Entonces:
\begin{align*}
\inner{\vect{x}}{\vect{y}} &= 1 \cdot 4 + 2 \cdot 5 + 3 \cdot 6 \\
&= 4 + 10 + 18 \\
&= 32
\end{align*}

\end{example}

\begin{definition}[Norma en $\R^n$]

Sea $\vect{x} \in \R^n$ con $\vect{x} = (x_1, x_2, \dots, x_n)$. La norma de $\vect{x}$ es el número real dado por:
\[
\norm{\vect{x}} = \sqrt{\inner{\vect{x}}{\vect{x}}}
\]

Expandiendo la definición, tenemos:
\begin{align*}
\norm{\vect{x}} &= \sqrt{x_1^2 + x_2^2 + \dots + x_n^2} \\
&= \sqrt{\sum_{i=1}^n x_i^2}
\end{align*}

Esta es conocida como norma inducida por el producto interior en $\R^n$.

\end{definition}

\begin{example}
  
Sea $\vect{x} = (3, 4)$. Entonces:
\begin{align*}
\norm{\vect{x}} &= \sqrt{3^2 + 4^2} \\
&= \sqrt{9 + 16} \\
&= \sqrt{25} \\
&= 5
\end{align*}

\end{example}

\section{Espacios en $\R^n$}

\begin{definition}[Bola abierta en $\R^n$]

Sea $\vect{a} \in \R^n$ y $r > 0$. La bola abierta de centro $\vect{a}$ y radio $r$ es el conjunto:
\[
B(\vect{a}, r) = \{ \vect{x} \in \R^n : \norm{\vect{x} - \vect{a}} < r \}
\]
\end{definition}

\begin{definition}[Bola cerrada en $\R^n$]

Sea $\vect{a} \in \R^n$ y $r > 0$. La bola cerrada de centro $\vect{a}$ y radio $r$ es el conjunto:
\[
\overline{B}(\vect{a}, r) = \{ \vect{x} \in \R^n : \norm{\vect{x} - \vect{a}} \leq r \}
\]
\end{definition}

\begin{definition}[Punto interior de un conjunto en $\R^n$]

Sea $A \subseteq \R^n$. Un punto $\vect{a} \in A$ es un punto interior de $A$ si y solo si:
\[
\exists r > 0 : B(\vect{a}, r) \subseteq A
\]
\end{definition}

\begin{definition}[Interior de un conjunto en $\R^n$]

Sea $A \subseteq \R^n$. El interior de $A$ es el conjunto de todos los puntos interiores de $A$, denotado por $\operatorname{Int}(A)$ o $A^\circ$:
\[
\operatorname{Int}(A) = \{ \vect{a} \in A : \vect{a} \text{ es un punto interior de } A \}
\]
\end{definition}

\begin{definition}[Punto frontera de un conjunto en $\R^n$]

Sea $A \subseteq \R^n$. Un punto $\vect{a} \in \R^n$ es un punto frontera de $A$ si y solo si:
\[
\forall r > 0 : (B(\vect{a}, r) \cap A \neq \emptyset) \wedge (B(\vect{a}, r) \cap (\R^n \setminus A) \neq \emptyset)
\]
\end{definition}

\begin{definition}[Frontera de un conjunto en $\R^n$]

Sea $A \subseteq \R^n$. La frontera de $A$ es el conjunto de todos los puntos frontera de $A$, denotado por $\partial A$:
\[
\partial A = \{ \vect{a} \in \R^n : \vect{a} \text{ es un punto frontera de } A \}
\]
\end{definition}

\begin{definition}[Punto adherente de un conjunto en $\R^n$]

Sea $A \subseteq \R^n$. Un punto $\vect{a} \in \R^n$ es un punto adherente de $A$ si y solo si:
\[
\forall r > 0 : B(\vect{a}, r) \cap A \neq \emptyset
\]
\end{definition}

\begin{definition}[Adherencia de un conjunto en $\R^n$]

Sea $A \subseteq \R^n$. La adherencia de $A$ es el conjunto de todos los puntos adherentes de $A$, denotado por $\overline{A}$:
\[
\overline{A} = \{ \vect{a} \in \R^n : \vect{a} \text{ es un punto adherente de } A \}
\]
\end{definition}

\begin{definition}[Punto de acumulación de un conjunto en $\R^n$]

Sea $A \subseteq \R^n$. Un punto $\vect{a} \in \R^n$ es un punto de acumulación de $A$ si y solo si:
\[
\forall r > 0 : (B(\vect{a}, r) \setminus \{\vect{a}\}) \cap A \neq \emptyset
\]
\end{definition}

\begin{definition}[Conjunto de acumulación de un conjunto en $\R^n$]

Sea $A \subseteq \R^n$. El conjunto de acumulación de $A$ es el conjunto de todos los puntos de acumulación de $A$, denotado por $A'$:
\[
A' = \{ \vect{a} \in \R^n : \vect{a} \text{ es un punto de acumulación de } A \}
\]
\end{definition}

\begin{proposition}
Sea $A \subseteq \R^n$. Entonces:

\begin{itemize}
\item $\operatorname{Int}(A) \cap \partial A = \emptyset$.
\item $\operatorname{Int}(A) \cup \partial A = \overline{A}$.
\item $\partial A \subseteq \overline{A}$.
\item $A' \subseteq \overline{A}$.
\end{itemize}
\end{proposition}

\begin{definition}[Conjunto cerrado en $\R^n$]

Sea $A \subseteq \R^n$. $A$ es un conjunto cerrado si y solo si $\partial A \subseteq A$.

\end{definition}

\begin{definition}[Conjunto abierto en $\R^n$]

Sea $A \subseteq \R^n$. $A$ es un conjunto abierto si y solo si $\partial A \cap A = \emptyset$.

\end{definition}

\begin{definition}[Conjunto acotado en $\R^n$]
Sea $A \subseteq \R^n$. $A$ es un conjunto acotado si y solo si:
\[
\exists r > 0 : A \subseteq B(\vect{0}, r)
\]
\end{definition}

\begin{definition}[Conjunto compacto en $\R^n$]
Sea $A \subseteq \R^n$. $A$ es un conjunto compacto si y solo si $A$ es cerrado y acotado.
\end{definition}

\begin{example}

Sea $A = \{ (x, y) \in \R^2 : x^2 + y^2 < 1 \}$. Entonces:

\begin{itemize}
\item $A^\circ = A$, pues cada punto de $A$ es un punto interior.
\item $\partial A = \{ (x, y) \in \R^2 : x^2 + y^2 = 1 \}$, pues cada punto de la circunferencia es un punto frontera.
\item $\overline{A} = A \cup \partial A$, pues cada punto de $A$ es un punto adherente y cada punto de $\partial A$ es un punto adherente.
\item $A' = A \cup \partial A$, pues cada punto de $A$ es un punto de acumulación y cada punto de $\partial A$ es un punto de acumulación.
\item $A$ no es cerrado, pues $\partial A \not\subseteq A$.
\item $A$ es abierto, pues $\partial A \cap A = \emptyset$.
\item $A$ es acotado, pues $A \subseteq B(\vect{0}, 2)$.
\item $A$ no es compacto, pues $A$ no es cerrado.
\end{itemize}

\end{example}

\section{Funciones de varias variables}

\begin{definition}[Función de varias variables]

Sea $f$ una función de $n$ variables, denotada por $f : U \subseteq \R^n \to \R$, es una regla de correspondencia que asigna a cada $n$-tupla $\vect{x} = (x_1, x_2, \dots, x_n) \in U$ un único $z \in \R$, denotado por:
\[
z = f(x_1, x_2, \dots, x_n)
\]
\end{definition}

\begin{definition}[Curvas de nivel]

Sea $f : U \subseteq \R^n \to \R$ una función de varias variables. La curva de nivel de $f$ asociada al valor $c \in \R$ es el conjunto:
\[
\{ \vect{x} \in U : f(\vect{x}) = c \}
\]
\end{definition}

\begin{example}
Sea $f : \R^2 \to \R$ dada por $f(x, y) = x^2 + y^2$. Entonces, la curva de nivel asociada al valor $c$ es el conjunto:
\[
\{ (x, y) \in \R^2 : x^2 + y^2 = c \}
\]
\end{example}

\end{document}